% ============================================================
%  pptheader.tex —— 主题、宏包、定理环境配置
%  在 template.tex 中通过 % ============================================================
%  pptheader.tex —— 主题、宏包、定理环境配置
%  在 template.tex 中通过 % ============================================================
%  pptheader.tex —— 主题、宏包、定理环境配置
%  在 template.tex 中通过 % ============================================================
%  pptheader.tex —— 主题、宏包、定理环境配置
%  在 template.tex 中通过 \input{pptheader.tex} 引入
% ============================================================

% ---- Beamer 主题与配色 ----
\usetheme{Madrid}
% 其他常用主题：Warsaw / Berlin / Singapore / AnnArbor / CambridgeUS

\usecolortheme{default}
% 其他常用配色：beaver / crane / orchid / seahorse

\usefonttheme{serif}       % 原 \documentclass[serif] 的正确写法

% 如需自定义主色调，取消下面两行注释并修改 RGB：
% \definecolor{themecolor}{RGB}{0, 82, 155}
% \usecolortheme[named=themecolor]{structure}

% ---- 字体 ----
\usepackage{fontspec}
\usepackage{xeCJK}
\setCJKmainfont{SimSun}   % 宋体；macOS/Linux 可换 STSong 或 FandolSong
\setCJKsansfont{SimHei}   % 黑体；macOS/Linux 可换 STHeiti 或 FandolHei

% ---- 数学 ----
\usepackage{amsmath}
\usepackage{amssymb}
\usepackage{amsthm}

% ---- 图表 ----
\usepackage{graphicx}
\graphicspath{{figures/}}   
\usepackage{subcaption}   % \begin{subfigure}
	\usepackage{float}        % 图表位置 [H]
	\usepackage{booktabs}     % 三线表 \toprule \midrule \bottomrule
	\usepackage{makecell}     % 表格内换行 \makecell{}
	
	% ---- 其他 ----
	% 注意：beamer 已内置 hyperref，无需重复引入
	
	% ---- 定理环境（配合 \documentclass[notheorems] 重新定义） ----
	\newtheorem{theorem}{定理}
	\newtheorem{lemma}{引理}
	\newtheorem{corollary}{推论}
	\newtheorem{definition}{定义}
	\newtheorem{remark}{注}
	
% ---- 个人简介页所需宏包 ----
\usepackage{tikz}
\usetikzlibrary{positioning}
\usetikzlibrary{arrows.meta}
\usepackage{pgfplots}         
\pgfplotsset{compat=1.18}      
\usepackage{tcolorbox}
\tcbuselibrary{skins}

% 从主题结构色自动派生，用 \AtBeginDocument 推迟到 Beamer 颜色初始化后执行
\AtBeginDocument{%
	\colorlet{panelbg}{structure.fg!12!white}   % 浅色面板背景
	\colorlet{tagborder}{structure.fg!60!white} % 标签边框
} 引入
% ============================================================

% ---- Beamer 主题与配色 ----
\usetheme{Madrid}
% 其他常用主题：Warsaw / Berlin / Singapore / AnnArbor / CambridgeUS

\usecolortheme{default}
% 其他常用配色：beaver / crane / orchid / seahorse

\usefonttheme{serif}       % 原 \documentclass[serif] 的正确写法

% 如需自定义主色调，取消下面两行注释并修改 RGB：
% \definecolor{themecolor}{RGB}{0, 82, 155}
% \usecolortheme[named=themecolor]{structure}

% ---- 字体 ----
\usepackage{fontspec}
\usepackage{xeCJK}
\setCJKmainfont{SimSun}   % 宋体；macOS/Linux 可换 STSong 或 FandolSong
\setCJKsansfont{SimHei}   % 黑体；macOS/Linux 可换 STHeiti 或 FandolHei

% ---- 数学 ----
\usepackage{amsmath}
\usepackage{amssymb}
\usepackage{amsthm}

% ---- 图表 ----
\usepackage{graphicx}
\graphicspath{{figures/}}   
\usepackage{subcaption}   % \begin{subfigure}
	\usepackage{float}        % 图表位置 [H]
	\usepackage{booktabs}     % 三线表 \toprule \midrule \bottomrule
	\usepackage{makecell}     % 表格内换行 \makecell{}
	
	% ---- 其他 ----
	% 注意：beamer 已内置 hyperref，无需重复引入
	
	% ---- 定理环境（配合 \documentclass[notheorems] 重新定义） ----
	\newtheorem{theorem}{定理}
	\newtheorem{lemma}{引理}
	\newtheorem{corollary}{推论}
	\newtheorem{definition}{定义}
	\newtheorem{remark}{注}
	
% ---- 个人简介页所需宏包 ----
\usepackage{tikz}
\usetikzlibrary{positioning}
\usetikzlibrary{arrows.meta}
\usepackage{pgfplots}         
\pgfplotsset{compat=1.18}      
\usepackage{tcolorbox}
\tcbuselibrary{skins}

% 从主题结构色自动派生，用 \AtBeginDocument 推迟到 Beamer 颜色初始化后执行
\AtBeginDocument{%
	\colorlet{panelbg}{structure.fg!12!white}   % 浅色面板背景
	\colorlet{tagborder}{structure.fg!60!white} % 标签边框
} 引入
% ============================================================

% ---- Beamer 主题与配色 ----
\usetheme{Madrid}
% 其他常用主题：Warsaw / Berlin / Singapore / AnnArbor / CambridgeUS

\usecolortheme{default}
% 其他常用配色：beaver / crane / orchid / seahorse

\usefonttheme{serif}       % 原 \documentclass[serif] 的正确写法

% 如需自定义主色调，取消下面两行注释并修改 RGB：
% \definecolor{themecolor}{RGB}{0, 82, 155}
% \usecolortheme[named=themecolor]{structure}

% ---- 字体 ----
\usepackage{fontspec}
\usepackage{xeCJK}
\setCJKmainfont{SimSun}   % 宋体；macOS/Linux 可换 STSong 或 FandolSong
\setCJKsansfont{SimHei}   % 黑体；macOS/Linux 可换 STHeiti 或 FandolHei

% ---- 数学 ----
\usepackage{amsmath}
\usepackage{amssymb}
\usepackage{amsthm}

% ---- 图表 ----
\usepackage{graphicx}
\graphicspath{{figures/}}   
\usepackage{subcaption}   % \begin{subfigure}
	\usepackage{float}        % 图表位置 [H]
	\usepackage{booktabs}     % 三线表 \toprule \midrule \bottomrule
	\usepackage{makecell}     % 表格内换行 \makecell{}
	
	% ---- 其他 ----
	% 注意：beamer 已内置 hyperref，无需重复引入
	
	% ---- 定理环境（配合 \documentclass[notheorems] 重新定义） ----
	\newtheorem{theorem}{定理}
	\newtheorem{lemma}{引理}
	\newtheorem{corollary}{推论}
	\newtheorem{definition}{定义}
	\newtheorem{remark}{注}
	
% ---- 个人简介页所需宏包 ----
\usepackage{tikz}
\usetikzlibrary{positioning}
\usetikzlibrary{arrows.meta}
\usepackage{pgfplots}         
\pgfplotsset{compat=1.18}      
\usepackage{tcolorbox}
\tcbuselibrary{skins}

% 从主题结构色自动派生，用 \AtBeginDocument 推迟到 Beamer 颜色初始化后执行
\AtBeginDocument{%
	\colorlet{panelbg}{structure.fg!12!white}   % 浅色面板背景
	\colorlet{tagborder}{structure.fg!60!white} % 标签边框
} 引入
% ============================================================

% ---- Beamer 主题与配色 ----
\usetheme{Madrid}
% 其他常用主题：Warsaw / Berlin / Singapore / AnnArbor / CambridgeUS

\usecolortheme{default}
% 其他常用配色：beaver / crane / orchid / seahorse

\usefonttheme{serif}       % 原 \documentclass[serif] 的正确写法

% 如需自定义主色调，取消下面两行注释并修改 RGB：
% \definecolor{themecolor}{RGB}{0, 82, 155}
% \usecolortheme[named=themecolor]{structure}

% ---- 字体 ----
\usepackage{fontspec}
\usepackage{xeCJK}
\setCJKmainfont{SimSun}   % 宋体；macOS/Linux 可换 STSong 或 FandolSong
\setCJKsansfont{SimHei}   % 黑体；macOS/Linux 可换 STHeiti 或 FandolHei

% ---- 数学 ----
\usepackage{amsmath}
\usepackage{amssymb}
\usepackage{amsthm}

% ---- 图表 ----
\usepackage{graphicx}
\graphicspath{{figures/}}   
\usepackage{subcaption}   % \begin{subfigure}
	\usepackage{float}        % 图表位置 [H]
	\usepackage{booktabs}     % 三线表 \toprule \midrule \bottomrule
	\usepackage{makecell}     % 表格内换行 \makecell{}
	
	% ---- 其他 ----
	% 注意：beamer 已内置 hyperref，无需重复引入
	
	% ---- 定理环境（配合 \documentclass[notheorems] 重新定义） ----
	\newtheorem{theorem}{定理}
	\newtheorem{lemma}{引理}
	\newtheorem{corollary}{推论}
	\newtheorem{definition}{定义}
	\newtheorem{remark}{注}
	
% ---- 个人简介页所需宏包 ----
\usepackage{tikz}
\usetikzlibrary{positioning}
\usetikzlibrary{arrows.meta}
\usepackage{pgfplots}         
\pgfplotsset{compat=1.18}      
\usepackage{tcolorbox}
\tcbuselibrary{skins}

% 从主题结构色自动派生，用 \AtBeginDocument 推迟到 Beamer 颜色初始化后执行
\AtBeginDocument{%
	\colorlet{panelbg}{structure.fg!12!white}   % 浅色面板背景
	\colorlet{tagborder}{structure.fg!60!white} % 标签边框
}